% SmartNotes – preamble for biologinotater (stil «Lærebok», farger fra «Salvie»).
% Claude får IKKE lov til å endre denne; alt innhold skrives i dokumentkroppen.
\documentclass[11pt,a4paper]{article}

\usepackage[T1]{fontenc}
\usepackage[utf8]{inputenc}
\usepackage[norsk]{babel}
\usepackage[sc]{mathpazo}
\linespread{1.06}
\usepackage{microtype}
\usepackage[a4paper,margin=2.3cm,headheight=14pt]{geometry}
\usepackage[parfill]{parskip}

% --- Matematikk ---
\usepackage{amsmath,amssymb,mathtools}
\usepackage{icomma}            % 9,81 skrives riktig i matematikkmodus
\usepackage{cancel}            % \cancel{} for forkorting
\usepackage{siunitx}
\sisetup{
  output-decimal-marker = {,},
  group-separator = {\,},
  group-minimum-digits = 5,
  exponent-product = \cdot,
  inter-unit-product = \cdot,
  per-mode = symbol,
  range-phrase = {\text{ til }},
  list-final-separator = {\text{ og }},
  list-pair-separator = {\text{ og }},
}
\usepackage[version=4]{mhchem} % prosesser: \ce{6 CO2 + 6 H2O -> C6H12O6 + 6 O2}

% Egne makroer (dokumentert i instruksene til Claude)
\DeclarePairedDelimiter{\abs}{\lvert}{\rvert}
\DeclarePairedDelimiter{\norm}{\lVert}{\rVert}
\newcommand{\dd}{\mathop{}\!\mathrm{d}}
\newcommand{\dv}[2]{\frac{\mathrm{d}#1}{\mathrm{d}#2}}
\newcommand{\ddv}[2]{\frac{\mathrm{d}^2#1}{\mathrm{d}#2^2}}
\newcommand{\pdv}[2]{\frac{\partial#1}{\partial#2}}
\newcommand{\vb}[1]{\mathbf{#1}}
\newcommand{\vu}[1]{\hat{\mathbf{#1}}}
\newcommand{\enhet}[1]{\,[\mathrm{#1}]}

% --- Figurer og tegninger ---
\usepackage{graphicx}
\usepackage{float}
\usepackage[labelfont=bf,font=small]{caption}
\usepackage{tikz}
\usetikzlibrary{arrows.meta,calc,angles,quotes,patterns,positioning,
  decorations.pathmorphing,decorations.markings,shapes.geometric,intersections,babel}
\usepackage{pgfplots}
\pgfplotsset{compat=1.18}
\usepackage[european,siunitx]{circuitikz}

% --- Tabeller og lister ---
\usepackage{booktabs,array,tabularx}
\usepackage{enumitem}

% --- Bokser ---
\usepackage{xcolor}
\usepackage[most]{tcolorbox}
% Salvie: dempet grønn og varme, rolige toner
\definecolor{snblue}{HTML}{3F5A6E}
\definecolor{sngreen}{HTML}{4A5D4A}
\definecolor{snorange}{HTML}{B07B3A}
\definecolor{sngray}{HTML}{6B6B66}
\definecolor{snyellow}{HTML}{9A7B2E}
\tcbset{snbase/.style={enhanced,breakable,frame hidden,boxrule=0pt,arc=0pt,
  left=9pt,right=6pt,top=5pt,bottom=5pt,fonttitle=\bfseries,
  detach title,before upper={\tcbtitle\par\smallskip},
  before skip=10pt,after skip=10pt}}

\newtcolorbox{definisjon}[1][]{snbase,colback=snblue!5,borderline west={2.4pt}{0pt}{snblue},coltitle=snblue!85!black,
  title={Definisjon\ifx\relax#1\relax\else: #1\fi}}
\newtcolorbox{formel}[1][]{snbase,colback=sngreen!5,borderline west={2.4pt}{0pt}{sngreen},coltitle=sngreen!85!black,
  title={\ifx\relax#1\relax Viktig\else #1\fi}}
\newtcolorbox{eksempel}[1][]{snbase,colback=sngray!5,borderline west={2.4pt}{0pt}{sngray},coltitle=sngray!85!black,
  title={Eksempel\ifx\relax#1\relax\else: #1\fi}}
\newtcolorbox{oppgave}[1][]{snbase,colback=snblue!5,borderline west={2.4pt}{0pt}{snblue},coltitle=snblue!85!black,
  title={Oppgave\ifx\relax#1\relax\else\ #1\fi}}
\newtcolorbox{losning}{snbase,colback=white,borderline west={1pt}{0pt}{snblue!50},coltitle=snblue!85!black,title={Løsning}}
\newtcolorbox{husk}{snbase,colback=snyellow!5,borderline west={2.4pt}{0pt}{snyellow},coltitle=snyellow!85!black,
  title={Husk}}
\newtcolorbox{merknad}{snbase,colback=snorange!5,borderline west={2.4pt}{0pt}{snorange},coltitle=snorange!85!black,
  title={Merknad fra SmartNotes}}

\newcommand{\uleselig}[1]{\textcolor{sngray}{\textit{#1}}\textsuperscript{\textcolor{snorange}{\,?}}}
\newcommand{\sidenotat}[1]{\marginpar{\footnotesize\raggedright #1}}

% Figurer klippet ut fra originalbildet. \notedir settes av SmartNotes.
\newcommand{\notedir}{.}
\newcommand{\originalfigur}[3][0.6]{%
  \begin{figure}[H]\centering
  \includegraphics[width=#1\linewidth,height=0.45\textheight,keepaspectratio]{\notedir/#2}%
  \ifx\relax#3\relax\else\caption{#3}\fi
  \end{figure}}
\newcommand{\manglerfigur}[1]{%
  \begin{center}\fbox{\parbox{0.6\linewidth}{\centering\small\textcolor{sngray}{[Figur mangler: #1]}}}\end{center}}

% --- Sidehode ---
\usepackage{fancyhdr}
\pagestyle{fancy}
\fancyhf{}
\renewcommand{\headrulewidth}{0.4pt}
\newcommand{\snfag}{}
\newcommand{\snkapittel}{}
\fancyhead[L]{\small\textcolor{sngray}{\snfag}}
\fancyhead[R]{\small\textcolor{sngray}{\snkapittel}}
\fancyfoot[C]{\small\thepage}

% Ingen nummerering av overskrifter i notater
\setcounter{secnumdepth}{0}
\setcounter{tocdepth}{1}

% Tittelblokk for ett notat
\newcommand{\notehode}[3]{%
  \begingroup
  \noindent{\LARGE\bfseries #1\par}\vspace{2pt}
  \noindent{\small\textcolor{sngray}{#2\ifx\relax#3\relax\else\ \textbullet\ #3\fi}\par}
  \vspace{4pt}\noindent\textcolor{snblue}{\rule{\linewidth}{1pt}}\par
  \endgroup\vspace{6pt}}

\usepackage{titlesec}
\titleformat{\section}{\Large\bfseries\color{sngreen!80!black}}{}{0pt}{}
\titleformat{\subsection}{\large\bfseries\color{sngreen!70!black}}{}{0pt}{}
\renewcommand{\notehode}[3]{%
  \begingroup
  \noindent{\Huge #1\par}\vspace{4pt}
  \noindent{\small\scshape\textcolor{sngray}{#2\ifx\relax#3\relax\else\ \textperiodcentered\ #3\fi}\par}
  \vspace{6pt}\noindent\textcolor{sngreen}{\rule{\linewidth}{0.6pt}}\par
  \endgroup\vspace{8pt}}

\usepackage[hidelinks]{hyperref}
\usepackage{bookmark}
